\documentclass[literature.tex]{subfiles}

\begin{document}
\prose{Spalovač mrtvol}{Ladislav Fuks}
\teorieA{
epická próza, psychologická novela (s prvky hororu, grotesky a černého humoru)
}{
Druhá vlna literatury o 2. světové válce a holokaustu; psychologická próza s existencialistickými a groteskními prvky.
}{
Spisovný jazyk s ironickým a chladným tónem. Časté opakování frází, zdrobnělin a přezdívek (Lakmé, Romane, čarokrásné). Odborné termíny z krematoria, nacistické fráze vytržené z kontextu, citace z knihy o Tibetu.
}{
Opakování (epizeuxis, anafora motivů), ironie, kontrast mezi něžným vyjadřováním a hrůznými činy.
}{
Alegorie (kremace = „spasení“ duší, plynové komory), symbolismus (bílá rakev, Tibet, popel), groteska (smrt jako „krásná věc“), oxymóron („šťastný nový řád“).
}{
Vyprávění v er-formě s vnitřními monology hlavní postavy. Opakující se motivy (růžolící dívka, šílená žena s pérem, novinové titulky) vytvářejí hypnotický, klaustrofobický dojem. Chladný, strohý popis hrůzných činů zesiluje morbiditu a psychologický horor.
}{
„Co kdybych tě drahá oběsil?“ Usmála se na něho dolů, snad mu dobře nerozuměla, on se usmál též, kopl do židle a bylo to.\\
„Smrt sbližuje,“ řekl si s rukou v kapse na kleštičkách, v lidském popelu není rozdíl. Jestli je to popel německého bannführera nebo čtvrtečního židovského chlapce, to je jedno.
}
\teorieB{
\wtem{Karel Kopfrkingl}{hlavní postava, zdánlivě jemný, tolerantní, abstinentní zaměstnanec krematoria, milující otec a manžel. Posedlý smrtí a kremací („spasení duší“). Lehce manipulovatelný, rychle se promění v udavače, fašistu a vraha vlastní rodiny. Nakonec zešílí a cítí se jako spasitel (Budha/dalajláma).}
\wtem{Marie}{hodná, milující žena poloviční Židovka, tolerantní k manželovým podivínstvím.}
\wtem{Zina a Milivoj}{dcera a syn, „čarokrásné“ a „slunné“ děti; Mili zženštilý, Zina upřímná. Oba zavražděni otcem.}
\wtem{Willy Reinke}{přítel z 1. světové války, nacistický fanatik (SdP/NSDAP), manipulátor, který probudí v Kopfrkinglovi „německou krev“.}
}{
Příběh sleduje postupnou proměnu Karla Kopfrkingla, zaměstnance pražského krematoria, v době před Mnichovem a na počátku okupace. Zpočátku je to spořádaný, mírumilovný podivín, který miluje rodinu, Tibet a kremaci jako „důstojné spasení“. Pod vlivem přítele Willyho Reinkeho (člena SdP) si uvědomí „německou krev“, začne věřit nacistické ideologii, vstoupí do strany, špehuje a udává Židy. Aby postoupil v kariéře (ředitel krematoria, později „experimentátor“ plynových komor), zavraždí ženu Lakmé (oběšení), syna Miliho (utlučen tyčí v krematoriu) a dcera Zina je rovněž odstraněna. Po vraždách zešílí, cítí se jako spasitel lidstva a nový dalajláma. Epilog po válce ukazuje jeho iluzi „nového šťastného řádu“.
}{
Krátká novela s chronologickým postupem a opakováním motivů (panoptikum, novinové fráze, krematorium). Prostor: Praha (krematorium, domov, synagoga). Čas: 1937–1939 (před Mnichovem a počátek okupace), krátký epilog po roce 1945.
}{
Varování před manipulovatelností „obyčejného“ člověka, extremismem a fašismem. Ukazuje, jak snadno se tolerantní maloměšťák promění v monstrum pod vlivem ideologie a kariérního postupu. Motivy: odcizení, banalita zla, posedlost smrtí, ztráta humanity („v lidském popelu není rozdíl“). Smrt jako „spasení“ předznamenává holokaust.
}\newpage
\historie{
Dílo napsáno v 60. letech (vydáno 1967), reflektuje předválečnou a okupační dobu (Mnichov 1938, nacistická okupace, antisemitismus). Fuks (sám Žid) zpracovává trauma holokaustu a mechanismy kolaborace.
}{}{}{}{}{
Ladislav Fuks (1923 Praha – 1994 Praha). Český prozaik židovského původu. Za války nucené práce, studoval filozofii a psychologii. Debutoval Mr. Theodor Mundstock (1963). Jeho tvorba často zpracovává téma holokaustu, úzkosti a manipulace psychiky. Spalovač mrtvol patří k jeho nejsilnějším dílům, zfilmováno Jurajem Herzem (1969, trezorový film).
}{
Mezi nejznámější díla patří \textsc{Mr. Theodor Mundstock}, \textsc{Variace pro temnou strunu}, \textsc{Mí černovlasí bratři}, \textsc{Příběh kriminálního rady}.
}
\kritika{}{}{
Téma manipulace ideologií, banality zla a ztráty vlastního názoru zůstává vysoce aktuální (populismus, extremismus, propaganda).
}
\nazor{
Spalovač mrtvol je mrazivě geniální studie, jak snadno se „slušný“ člověk promění v monstrum. Kopfrkingl zpočátku působí jako neškodný podivín, ale pod vlivem Willyho a kariérního lákadla se stává vrahem vlastní rodiny – a přitom si stále myslí, že je „jemný“ a spasitel. Opakující se motivy a chladný popis vražd vyvolávají nepříjemné mrazení. Fuks mistrně ukazuje banalitu zla. Kniha je krátká, ale intenzivní – rozhodně jedna z nejsilnějších psychologických novel o fašismu v české literatuře. Doporučuji, ale s opatrností (je to dost drsné).
}
\explain{
\textbl{Groteska --} spojení komického a hrůzného, často absurdního;
\textbl{Banalita zla --} zdánlivě obyčejný člověk páchá zrůdnosti;
\textbl{Psychologický horor --} zaměření na vnitřní proměnu a úzkost postavy.
}
\wpage
\end{document}